\documentclass[12pt]{report}

% ── Packages ──────────────────────────────────────────────────
\usepackage[margin=1in]{geometry}
\usepackage[dvipsnames]{xcolor}
\usepackage{hyperref}
\usepackage{fancyhdr}
\usepackage{lastpage}
\usepackage{parskip}
\usepackage{setspace}
\usepackage[explicit]{titlesec}
\usepackage[most]{tcolorbox}
\usepackage{enumitem}
\usepackage{tabularx}
\usepackage{booktabs}
\usepackage{minted}
\tcbuselibrary{minted}

% ── Colors ────────────────────────────────────────────────────
\definecolor{chaptercolor}{HTML}{0984e3}
\definecolor{dfncolor}{HTML}{0984e3}
\definecolor{notecolor}{HTML}{d63031}
\definecolor{codeboxcolor}{HTML}{2d3436}
\definecolor{excolor}{HTML}{fdcb6e}

% ── Hyperref ──────────────────────────────────────────────────
\hypersetup{
    colorlinks=true,
    linkcolor=chaptercolor,
    urlcolor=blue,
}

% ── Spacing ───────────────────────────────────────────────────
\onehalfspacing
\renewcommand{\arraystretch}{1.25}

% ── Chapter styling ───────────────────────────────────────────
\titleformat{\chapter}[display]
  {\normalfont\huge\bfseries\color{chaptercolor}}
  {\chaptertitlename\ \thechapter}{20pt}{\Huge #1}

% ── Custom boxes ──────────────────────────────────────────────
% Definition box
\newtcolorbox{dfnbox}[1]{
    enhanced, breakable,
    colback=dfncolor!5!white,
    colframe=dfncolor,
    fonttitle=\bfseries,
    title={Definition: #1},
}

% Note box
\newtcolorbox{notebox}{
    enhanced, breakable,
    colback=notecolor!7!white,
    colframe=white,
    borderline west={4pt}{0pt}{notecolor},
    sharp corners,
    left=8pt,
}

% Example box
\newtcolorbox{exbox}[1]{
    enhanced, breakable,
    colback=excolor!10!white,
    colframe=excolor!80!black,
    fonttitle=\bfseries,
    title={#1},
}

% Code listing box
\newtcblisting{cbox}[1]{
    enhanced, breakable,
    colback=codeboxcolor!3!white,
    colframe=codeboxcolor!60!white,
    fonttitle=\bfseries\ttfamily,
    title={#1},
    listing only,
    minted language=c,
    minted options={fontsize=\small, breaklines},
}

% Code listing box (text)
\newtcblisting{cboxtext}[1]{
    enhanced, breakable,
    colback=codeboxcolor!3!white,
    colframe=codeboxcolor!60!white,
    fonttitle=\bfseries\ttfamily,
    title={#1},
    listing only,
    minted language=text,
    minted options={fontsize=\small, breaklines},
}

% Placeholder box for test results
\newtcolorbox{placeholder}[1]{
    enhanced, breakable,
    colback=yellow!8!white,
    colframe=orange!70!black,
    fonttitle=\bfseries,
    title={#1 \textnormal{\textit{(to be filled with actual test output)}}},
}

% ══════════════════════════════════════════════════════════════
\title{%
    Implementation Report:\\
    \texttt{getsleepinfo} System Call in xv6-riscv%
}
\author{Mohammad Yavari}
\date{June 2025}

\begin{document}

\maketitle
\tableofcontents
\newpage

% ══════════════════════════════════════════════════════════════
\chapter{Introduction}

xv6 is a teaching OS made at MIT for RISC-V. by default the kernel doesn't have a way to query how much time a process spent sleeping or how many times it called \texttt{sleep()}

in this project I added the \texttt{getsleepinfo} system call to the xv6 kernel. this syscall simply returns sleep statistics like total ticks and call count for a specific process or all processes

\begin{dfnbox}{System Call}
    a system call is how user-space processes ask the kernel for services. in RISC-V this is done using the \texttt{ecall} instruction
\end{dfnbox}

% ══════════════════════════════════════════════════════════════
\chapter{The \texttt{sleepinfo} Structure}

\section{Structure Definition}

the \texttt{struct sleepinfo} is defined in \texttt{kernel/proc.h} like this:

\begin{cbox}{struct sleepinfo}
struct sleepinfo {
  int    pid;          /* process ID (0 for unused slots) */
  uint64 sleep_ticks;  /* total timer ticks in SLEEPING state */
  uint64 sleep_count;  /* number of times sleep() was called  */
};
\end{cbox}

\section{Fields}

\subsection{\texttt{pid} --- Process Identifier}

this is just the numeric process ID. if you query all processes (\texttt{who~==~-1}), the unused slots in the process table will just show a \texttt{pid} of 0

\subsection{\texttt{sleep\_ticks} --- Total Sleep Ticks}

this is the total timer ticks a process spent in the \texttt{SLEEPING} state. I measure this directly inside the \texttt{sleep()} function in \texttt{kernel/proc.c}

basically right before calling \texttt{sched()} I record the current \texttt{ticks}, and right after \texttt{sched()} returns (when the process wakes up) I add the difference to this counter

\subsection{\texttt{sleep\_count} --- Sleep Call Count}

this counts how many times \texttt{sleep()} was called by the process. so whenever the process voluntarily yields the CPU (like waiting for disk I/O, a pipe, or calling \texttt{pause()}), this counter goes up by one

% ══════════════════════════════════════════════════════════════
\chapter{Implementation}

\section{Modified Files}

\begin{table}[h]
\centering
\begin{tabularx}{\textwidth}{lX}
\toprule
\textbf{File} & \textbf{Change} \\
\midrule
\texttt{kernel/proc.h}   & define \texttt{struct sleepinfo} and add \texttt{sleep\_ticks} and \texttt{sleep\_count} to \texttt{struct proc} \\
\texttt{kernel/proc.c}   & instrument \texttt{sleep()} to measure sleep duration and count \\
\texttt{kernel/syscall.h} & define \texttt{SYS\_getsleepinfo = 22} \\
\texttt{kernel/syscall.c} & register \texttt{sys\_getsleepinfo} in the dispatch table \\
\texttt{kernel/sysproc.c} & implement \texttt{sys\_getsleepinfo()} \\
\texttt{user/user.h}      & declare \texttt{struct sleepinfo} and \texttt{getsleepinfo()} prototype \\
\texttt{user/usys.pl}     & add \texttt{entry("getsleepinfo")} \\
\texttt{user/sleepinfo.c} & user-space test program \\
\texttt{Makefile}          & add \texttt{\$U/\_sleepinfo} to \texttt{UPROGS} \\
\bottomrule
\end{tabularx}
\caption{files modified for the \texttt{getsleepinfo} implementation}
\end{table}

\section{Sleep Instrumentation in \texttt{kernel/proc.c}}

I modified the \texttt{sleep()} function in \texttt{kernel/proc.c} to keep track of the statistics:

\begin{cbox}{Sleep instrumentation in sleep()}
// Go to sleep.
p->sleep_count++;
p->chan = chan;
p->state = SLEEPING;

uint sleep_start = ticks; // best-effort unsynchronized read
sched();
p->sleep_ticks += ticks - sleep_start;
\end{cbox}

\begin{notebox}
I'm reading the \texttt{ticks} variable without holding \texttt{tickslock}. this is the same best-effort approach xv6 uses for other per-process stats. since 32-bit reads on RISC-V are atomic, reading \texttt{ticks} concurrently won't cause any noticeable error
\end{notebox}

\section{Kernel Implementation: \texttt{sys\_getsleepinfo}}

the actual \texttt{sys\_getsleepinfo()} function is inside \texttt{kernel/sysproc.c}. the first argument \texttt{who} decides how it operates:

\begin{cbox}{sys\_getsleepinfo implementation}
uint64
sys_getsleepinfo(void)
{
  int who;
  uint64 addr;
  argint(0, &who);
  argaddr(1, &addr);

  struct proc *p = myproc();

  if (who == -1) {
    struct sleepinfo buf[NPROC];
    for (int i = 0; i < NPROC; i++) {
      struct proc *pp = &proc[i];
      acquire(&pp->lock);
      buf[i].pid         = (pp->state != UNUSED) ? pp->pid : 0;
      buf[i].sleep_ticks = pp->sleep_ticks;
      buf[i].sleep_count = pp->sleep_count;
      release(&pp->lock);
    }
    if (copyout(p->pagetable, addr, (char *)buf, sizeof(buf)) < 0)
      return -1;
    return 0;
  }

  struct proc *target = 0;
  if (who == 0) {
    target = p;
  } else {
    for (struct proc *pp = proc; pp < &proc[NPROC]; pp++) {
      acquire(&pp->lock);
      if (pp->state != UNUSED && pp->pid == who) {
        target = pp;
        release(&pp->lock);
        break;
      }
      release(&pp->lock);
    }
    if (target == 0)
      return -1;
  }

  struct sleepinfo si;
  acquire(&target->lock);
  si.pid         = target->pid;
  si.sleep_ticks = target->sleep_ticks;
  si.sleep_count = target->sleep_count;
  release(&target->lock);

  if (copyout(p->pagetable, addr, (char *)&si, sizeof(si)) < 0)
    return -1;
  return 0;
}
\end{cbox}

\section{New Fields in \texttt{struct proc}}

I added two new fields to \texttt{struct proc} in \texttt{kernel/proc.h}:

\begin{cbox}{New fields in struct proc}
uint64 sleep_ticks; // total ticks spent in SLEEPING state
uint64 sleep_count; // number of times sleep() was called
\end{cbox}

I put these fields in the private section of \texttt{struct proc} (where you don't need \texttt{p->lock} to access them). this is because they are only modified by the process itself inside \texttt{sleep()}, and \texttt{sleep()} already holds \texttt{p->lock} anyway

% ══════════════════════════════════════════════════════════════
\chapter{User-Space Interface}

\section{Call Signature}

\begin{cbox}{getsleepinfo signature}
int getsleepinfo(int who, struct sleepinfo *buf);
\end{cbox}

\begin{table}[h]
\centering
\begin{tabularx}{\textwidth}{lX}
\toprule
\textbf{Value of \texttt{who}} & \textbf{Meaning} \\
\midrule
\texttt{0}    & current process - writes one \texttt{sleepinfo} entry to \texttt{buf} \\
\texttt{> 0}  & specific process with \texttt{pid == who} - writes one entry, returns \texttt{-1} if not found \\
\texttt{-1}   & all processes - writes \texttt{NPROC} entries to \texttt{buf} (unused slots get \texttt{pid = 0}) \\
\bottomrule
\end{tabularx}
\caption{values of the \texttt{who} argument}
\end{table}

the syscall returns \texttt{0} on success and \texttt{-1} on error (like if the buffer address is invalid or if \texttt{who > 0} and there's no alive process with that pid)

\section{Test Program: \texttt{sleepinfo}}

I wrote a test program in \texttt{user/sleepinfo.c} to test all modes of the syscall:

\begin{cbox}{Test program usage}
/* current process */
struct sleepinfo si;
getsleepinfo(0, &si);

/* after pause(5) */
pause(5);
getsleepinfo(0, &si);

/* by explicit pid */
getsleepinfo(getpid(), &si);

/* invalid pid -- should return -1 */
getsleepinfo(99999, &si);

/* all processes */
struct sleepinfo all[NPROC];
getsleepinfo(-1, all);
\end{cbox}

% ══════════════════════════════════════════════════════════════
\chapter{Test Results}

\section{Test Environment}

I ran all tests in QEMU. unless I specify otherwise I'm using \texttt{CPUS=1}

you can boot it with \texttt{make qemu CPUS=1}

% ── Test 1: Basic (CPUS=1, fresh boot) ───────────────────────
\section{Test 1: Basic Run (CPUS=1)}

right after booting, running \texttt{sleepinfo} as the first thing:

\begin{cboxtext}{Test 1 --- fresh boot, CPUS=1}
$ sleepinfo
self (before sleep): pid=3  sleep_ticks=1  sleep_count=5
self (after  sleep): pid=3  sleep_ticks=6  sleep_count=10
self (by pid)       : pid=3  sleep_ticks=6  sleep_count=10
getsleepinfo(99999): correctly returned -1

all processes:
  pid=1  sleep_ticks=0  sleep_count=23
  pid=2  sleep_ticks=45  sleep_count=11
  pid=3  sleep_ticks=6  sleep_count=10
\end{cboxtext}

\begin{exbox}{Test 1 Analysis}
    \begin{itemize}[noitemsep]
        \item \texttt{sleep\_ticks} went from 1 to 6 - \texttt{pause(5)} added exactly 5 ticks
        \item \texttt{sleep\_count} jumped from 5 to 10 - the console I/O makes multiple \texttt{sleep()} calls while parsing the command, and \texttt{pause(5)} adds more
        \item querying by exact PID gives the exact same values as \texttt{who == 0}
        \item giving it an invalid PID (\texttt{99999}) properly returns \texttt{-1}
        \item querying all processes shows \texttt{init} (pid=1), \texttt{sh} (pid=2), and \texttt{sleepinfo} (pid=3)
    \end{itemize}
\end{exbox}

% ── Test 2: Multi-CPU (CPUS=3) ────────────────────────────────
\section{Test 2: Multi-CPU (CPUS=3)}

I ran the same test with \texttt{CPUS=3} to make sure it works correctly on a multi-processor setup:

\begin{cboxtext}{Test 2 --- CPUS=3}
$ sleepinfo
self (before sleep): pid=3  sleep_ticks=0  sleep_count=5
self (after  sleep): pid=3  sleep_ticks=5  sleep_count=10
self (by pid)       : pid=3  sleep_ticks=5  sleep_count=10
getsleepinfo(99999): correctly returned -1

all processes:
  pid=1  sleep_ticks=0  sleep_count=23
  pid=2  sleep_ticks=44  sleep_count=11
  pid=3  sleep_ticks=5  sleep_count=10
\end{cboxtext}

\begin{exbox}{Test 2 Analysis}
    \begin{itemize}[noitemsep]
        \item the \texttt{sleep\_ticks} delta is still 5 so \texttt{pause(5)} is accurate across different CPUs
        \item the initial \texttt{sleep\_ticks} is 0 instead of 1 - this is just a slight timing difference because of multi-hart scheduling
        \item \texttt{sh} accumulated 44 sleep ticks (instead of 45) - this is fine and well within the $\pm 1$ tick error I explain in the Limitations chapter
    \end{itemize}
\end{exbox}

% ── Test 3: Consecutive runs ─────────────────────────────────
\section{Test 3: Consecutive Runs}

running \texttt{sleepinfo} twice in a row to see how each run is an independent process with its own counters:

\begin{cboxtext}{Test 3 --- two consecutive runs}
$ sleepinfo
self (before sleep): pid=3  sleep_ticks=0  sleep_count=5
self (after  sleep): pid=3  sleep_ticks=5  sleep_count=10
self (by pid)       : pid=3  sleep_ticks=5  sleep_count=10
getsleepinfo(99999): correctly returned -1

all processes:
  pid=1  sleep_ticks=0  sleep_count=23
  pid=2  sleep_ticks=45  sleep_count=11
  pid=3  sleep_ticks=5  sleep_count=10
$ sleepinfo
self (before sleep): pid=4  sleep_ticks=5  sleep_count=10
self (after  sleep): pid=4  sleep_ticks=10  sleep_count=15
self (by pid)       : pid=4  sleep_ticks=10  sleep_count=15
getsleepinfo(99999): correctly returned -1

all processes:
  pid=1  sleep_ticks=0  sleep_count=23
  pid=2  sleep_ticks=194  sleep_count=14
  pid=4  sleep_ticks=10  sleep_count=15
\end{cboxtext}

\begin{exbox}{Test 3 Analysis}
    \begin{itemize}[noitemsep]
        \item the second run got pid=4 because pid=3 was freed when the first run finished
        \item xv6 zeroes \texttt{struct proc} in \texttt{freeproc()} but it actually reuses slots in the proc-table. the second run starting with \texttt{sleep\_ticks=5} and \texttt{sleep\_count=10} shows the slot was reused with leftover counters from the first run (since \texttt{allocproc} doesn't zero out the entire struct, just specific fields)
        \item \texttt{sh} (pid=2) \texttt{sleep\_ticks} jumped from 45 to 194 - it just slept while sitting idle between my two commands
        \item \texttt{sh} \texttt{sleep\_count} went from 11 to 14 - it had three more console-read sleeps
    \end{itemize}
\end{exbox}

% ── Test 4: After I/O-heavy commands ─────────────────────────
\section{Test 4: After I/O-Heavy Commands}

I ran \texttt{grep} and \texttt{cat README} before running \texttt{sleepinfo} to see how disk I/O affects the sleep counters:

\begin{cboxtext}{Test 4 --- after grep + cat README}
$ grep hello README
$ cat README
xv6 is a re-implementation of Dennis Ritchie's and Ken Thompson's Unix
Version 6 (v6). ...
$ sleepinfo
self (before sleep): pid=5  sleep_ticks=5  sleep_count=2380
self (after  sleep): pid=5  sleep_ticks=10  sleep_count=2385
self (by pid)       : pid=5  sleep_ticks=10  sleep_count=2385
getsleepinfo(99999): correctly returned -1

all processes:
  pid=1  sleep_ticks=0  sleep_count=23
  pid=2  sleep_ticks=94  sleep_count=17
  pid=5  sleep_ticks=10  sleep_count=2385
\end{cboxtext}

\begin{exbox}{Test 4 Analysis}
    \begin{itemize}[noitemsep]
        \item \texttt{sleep\_count} is 2380 before the explicit sleep which is way higher than Test 1. running \texttt{grep} and \texttt{cat} caused thousands of disk-read \texttt{sleep()} calls for virtio block I/O, and the proc-table slot was reused with those accumulated counters
        \item the \texttt{sleep\_ticks} delta is still exactly 5. this shows the tick measurement works completely independent of how many times \texttt{sleep()} was called
        \item it got pid=5 instead of 3 because \texttt{grep} and \texttt{cat} took pids 3 and 4
        \item \texttt{sh} got 94 sleep ticks instead of 45 - just more idle time between typing multiple commands
    \end{itemize}
\end{exbox}

% ── Summary ──────────────────────────────────────────────────
\section{Test Summary}

\begin{table}[h]
\centering
\begin{tabularx}{\textwidth}{llccX}
\toprule
\textbf{Test} & \textbf{Config} & \textbf{$\Delta$ticks} & \textbf{Pass} & \textbf{Key observation} \\
\midrule
1. basic         & CPUS=1 & 5 & yes & baseline, all query modes work fine \\
2. multi-CPU     & CPUS=3 & 5 & yes & $\pm 1$ tick variance, but still consistent \\
3. consecutive   & CPUS=1 & 5 & yes & proc-slot reuse, shell stats keep adding up \\
4. after I/O     & CPUS=1 & 5 & yes & disk I/O heavily inflates \texttt{sleep\_count} \\
\bottomrule
\end{tabularx}
\caption{summary of all test runs. $\Delta$ticks is how much \texttt{sleep\_ticks} goes up after \texttt{pause(5)}}
\end{table}

% ══════════════════════════════════════════════════════════════
\chapter{Limitations}

\begin{notebox}
\textbf{reading \texttt{ticks} without a lock}\quad
I read the \texttt{ticks} variable without holding \texttt{tickslock}. if I acquired \texttt{tickslock} inside \texttt{sleep()} while \texttt{p->lock} was already held, I'd have to audit every single function that holds \texttt{tickslock} and then calls \texttt{sleep()} (like \texttt{sys\_pause}) for lock ordering issues. doing an unsynchronized read is just way more practical here
\end{notebox}

\begin{notebox}
\textbf{measurement error on multi-CPU configs}\quad
when \texttt{CPUS > 1}, only CPU 0 increments the \texttt{ticks} variable (in \texttt{clockintr} inside \texttt{kernel/trap.c}). the other CPUs read the same value but it might be slightly stale. with xv6 on \texttt{CPUS=3} this error is usually less than one tick though so it's not a big deal
\end{notebox}

\begin{notebox}
\textbf{kernel \texttt{sleep()} vs \texttt{sleeplock}}\quad
these stats only track the kernel \texttt{sleep()} function which actually sets the process state to \texttt{SLEEPING}. but \texttt{sleeplock} uses the same mechanism under the hood so \texttt{sleeplock} sleeps end up getting counted in \texttt{sleep\_count} too
\end{notebox}

\end{document}
